\chapter{Worked Makefiles}\label{ch:mk-worked}

Three Makefiles worth stealing. Each uses only what Part~\ref{part:make} has
covered, and each is short enough to read in one sitting.

\section{A Multi-File C Project}\label{sec:mk-cproject}

Sources in \file{src/}, headers in \file{include/}, everything generated under
\file{build/}, header dependencies maintained by the compiler.

\begin{makecode}
# ---- configuration (override any of these on the command line) -----------
CC       ?= gcc
CPPFLAGS ?= -Iinclude -MMD -MP
CFLAGS   ?= -Wall -Wextra -Wpedantic -std=c17 -O2 -g
LDFLAGS  ?=
LDLIBS   ?= -lm
PREFIX   ?= /usr/local

TARGET   := app
BUILD    := build
SRC      := $(wildcard src/*.c)
OBJ      := $(patsubst src/%.c,$(BUILD)/%.o,$(SRC))
DEP      := $(OBJ:.o=.d)

# ---- build ---------------------------------------------------------------
all: $(TARGET)

$(TARGET): $(OBJ)
    $(CC) $(LDFLAGS) -o $@ $^ $(LDLIBS)

$(BUILD)/%.o: src/%.c | $(BUILD)
    $(CC) $(CPPFLAGS) $(CFLAGS) -c -o $@ $<

$(BUILD):
    mkdir -p $@

# ---- variants ------------------------------------------------------------
debug: CFLAGS := -Wall -Wextra -std=c17 -O0 -g -fsanitize=address,undefined
debug: LDFLAGS += -fsanitize=address,undefined
debug: clean all

# ---- housekeeping --------------------------------------------------------
run: $(TARGET)
    ./$(TARGET)

test: $(TARGET)
    ./tests/run.sh

install: $(TARGET)
    install -d $(DESTDIR)$(PREFIX)/bin
    install -m 755 $(TARGET) $(DESTDIR)$(PREFIX)/bin/$(TARGET)

clean:
    rm -rf $(BUILD) $(TARGET)

print-%:
    @echo '$* = [$($*)]'

-include $(DEP)

.PHONY: all debug run test install clean
\end{makecode}

\begin{note}
	Every idea in Part~\ref{part:make} is here: \texttt{?=} for configurability
	(\autoref{sec:mk-condassign}), automatic variables
	(\autoref{sec:mk-auto}), a pattern rule (\autoref{sec:mk-pattern}), an
	order-only prerequisite for the output directory
	(\autoref{sec:mk-orderonly}), generated dependencies
	(\autoref{sec:mk-depgen}) and \texttt{.PHONY}
	(\autoref{sec:mk-phonyt}). It works unchanged for one source file or fifty,
	and it is correct under \cmd{make -j}.
\end{note}

\begin{note}
	\texttt{debug} depends on \texttt{clean} deliberately. As
	\autoref{sec:mk-override} warned, changing a flag changes no timestamp, so
	without the \cmd{clean} you would link sanitiser-instrumented objects against
	ordinary ones. The alternative --- a separate output directory per variant, as
	in \autoref{sec:mk-static} --- avoids the rebuild but costs more Makefile.
\end{note}

\section{A \texorpdfstring{\LaTeX{}}{LaTeX} Document}\label{sec:mk-latex}

\LaTeX{} is a bad fit for make's model. The number of runs needed is not known
in advance: a document with cross-references needs at least two passes, three if
the table of contents shifts a page, plus \cmd{biber} somewhere in the middle.
Make cannot express ``repeat until the \file{.aux} file stops changing''.

\cmd{latexmk} can, so the honest Makefile delegates to it:

\begin{makecode}
MAIN  := CSTools
ENGINE := -xelatex
FLAGS := -interaction=nonstopmode -synctex=1 -output-directory=build

all: $(MAIN).pdf

$(MAIN).pdf: $(wildcard *.tex) $(wildcard Chapters/*.tex) reference.bib
    latexmk $(ENGINE) $(FLAGS) $(MAIN).tex

watch:
    latexmk $(ENGINE) $(FLAGS) -pvc $(MAIN).tex

clean:
    latexmk -C -output-directory=build $(MAIN).tex

.PHONY: all watch clean
\end{makecode}

\begin{note}
	The prerequisites are still worth listing, even though \cmd{latexmk} does its
	own dependency tracking: they let \cmd{make} skip the whole thing when nothing
	has changed, and they make \cmd{make -n} honest about whether a build is
	needed.
\end{note}

\begin{moral}
	When a tool already solves the recurrence problem --- \cmd{latexmk},
	\cmd{cargo}, \cmd{go build} --- make's job is to call it, not to reimplement
	it. A Makefile that runs \cmd{xelatex} three times and hopes is worse than the
	one line it replaced.
\end{moral}

\begin{remark}
	The one place a Makefile earns its keep in a \LaTeX{} project is generated
	figures --- diagrams compiled from source, plots produced by a script, images
	converted between formats. Those \emph{are} a timestamp problem, and
	\autoref{sec:mk-pipeline} is the shape to use.
\end{remark}

\section{A Small Data Pipeline}\label{sec:mk-pipeline}

Nothing about make is specific to compilers. Any process that turns files into
other files fits, and a pipeline is where the incremental rebuild is most
valuable --- the steps are slow and you will run them repeatedly.

\begin{makecode}
RAW   := $(wildcard data/raw/*.csv)
CLEAN := $(patsubst data/raw/%.csv,data/clean/%.csv,$(RAW))
FIGS  := $(patsubst data/clean/%.csv,figures/%.png,$(CLEAN))

all: report.pdf

# stage 1: clean each raw file
data/clean/%.csv: data/raw/%.csv scripts/clean.py | data/clean
    python3 scripts/clean.py $< $@

# stage 2: one figure per cleaned file
figures/%.png: data/clean/%.csv scripts/plot.py | figures
    python3 scripts/plot.py $< $@

# stage 3: everything feeds the report
report.pdf: report.tex $(FIGS) summary.csv
    latexmk -xelatex -interaction=nonstopmode report.tex

summary.csv: $(CLEAN) scripts/summarise.py
    python3 scripts/summarise.py $(CLEAN) > $@

data/clean figures:
    mkdir -p $@

clean:
    rm -rf data/clean figures summary.csv report.pdf

.PHONY: all clean
\end{makecode}

\begin{note}
	Notice that \file{scripts/clean.py} is a prerequisite of the files it produces.
	Editing the script rebuilds everything downstream of it, which is exactly
	right and is the step people leave out. Without that edge, fixing a bug in the
	cleaning code leaves the old, wrong output in place and make reports success.
\end{note}

\begin{eg}
	Adding one raw file to \file{data/raw/} and running \cmd{make} cleans it,
	plots it, regenerates \file{summary.csv} --- because the summary depends on all
	the cleaned files --- and rebuilds the report. Nothing else runs. With
	\cmd{make -j4}, the per-file stages overlap.
\end{eg}

\begin{moral}
	If you find yourself keeping a note of which script to re-run after which
	change, you have a dependency graph in your head. Write it down as a Makefile
	and stop remembering it.
\end{moral}

\begin{tryit}
	Take any multi-step process you currently run by hand --- a build, a document,
	a data set --- and write the graph before writing any rules. Name the files at
	each stage and draw the arrows between them, as on
	\autopageref{sec:mk-graph}. The Makefile is then a transcription, and it will
	take about ten minutes.
\end{tryit}
